%% 7-0 序章　表のままで計算できたら？（第14回：組 V7-1。ページ制限なし。偶数ページ）
%% ① ゲームのキャラクターが回る（はてな？）② 方眼の家の絵を90°回す（やってみよう、4-8 の (x,y)→(-y,x)）
%% ③ 4つの数の表で絵を変える（主役の行列の紹介）④ 九九の表（1-13 の約束の回収）⑤ 表の路線 ⑥ 行列が働く場面
%% ⑦ この巻の地図と主役の行列の6つの顔 ⑧ 表計算と Python ⑨ 新しい記号 ⑩ この巻の読み方
%% 型（家の絵・行と列をなぞる図・表計算・Python）は tools/kata7.md。家の頂点 (0,0)(2,0)(2,2)(1,3)(0,2)、戸は右寄り。
\usetikzlibrary{matrix}
\makeatletter
\gdef\bk@partlabel{第7巻}\gdef\bk@partname{まとめて計算する}
%% ---- この章だけで使う部品（頭文字 \kviia…） ----
\tikzset{
  kviia mat/.style={matrix of math nodes, ampersand replacement=\&,
    left delimiter=(, right delimiter=),
    nodes={minimum width=7mm, minimum height=6mm, inner sep=0pt, anchor=center},
    column sep=0pt, row sep=0pt, inner sep=0.6mm},
  kviia lab/.style={font=\sffamily\footnotesize, inner sep=1pt}}
% 方眼と軸（tikzpicture[x=5mm, y=5mm] の中で）  \kviiagrid{xmin}{xmax}{ymin}{ymax}
\newcommand\kviiagrid[4]{%
  \draw[bkLine, line width=0.4pt, step=1] (#1,#3) grid (#2,#4);
  \draw[bkInk, line width=0.7pt, -{Stealth[length=2.2mm]}] (#1,0) -- ({#2+0.7},0) node[right, font=\footnotesize, inner sep=1pt] {$x$};
  \draw[bkInk, line width=0.7pt, -{Stealth[length=2.2mm]}] (0,#3) -- (0,{#4+0.7}) node[above, font=\footnotesize, inner sep=1pt] {$y$};
  \node[below left, font=\sffamily\footnotesize, inner sep=1.2pt] at (0,0) {O};}
% もとの家（灰色の点線）
\newcommand\kviiahousezero{%
  \filldraw[fill=bkGray!12, draw=bkGray, line width=0.8pt, dashed, line join=round] (0,0) -- (2,0) -- (2,2) -- (1,3) -- (0,2) -- cycle;
  \draw[bkGray, line width=0.6pt, dashed] (1.25,0) -- (1.25,1) -- (1.75,1) -- (1.75,0);}
% 行列 (a b; c d) をかけた家  \kviiahouse[色]{a}{b}{c}{d}
\newcommand\kviiahouse[5][bkOrange]{%
  \begin{scope}[cm={#2,#4,#3,#5,(0,0)}]
    \filldraw[fill=#1!25, draw=#1!80!black, line width=1.1pt, line join=round] (0,0) -- (2,0) -- (2,2) -- (1,3) -- (0,2) -- cycle;
    \filldraw[fill=white, draw=#1!80!black, line width=0.8pt] (1.25,0) -- (1.75,0) -- (1.75,1) -- (1.25,1) -- cycle;
  \end{scope}}
% ゲームのキャラクター（小さなロボット）  \kviiachara{(x,y)}{角度}{色}
\newcommand\kviiachara[3]{\begin{scope}[shift={#1}, rotate=#2]
  \filldraw[fill=#3!30, draw=#3!80!black, line width=0.9pt, rounded corners=1mm] (-5,-6) rectangle (5,4);
  \filldraw[fill=white, draw=#3!80!black, line width=0.8pt, rounded corners=0.8mm] (-4,4.6) rectangle (4,11);
  \fill[bkInk] (-1.6,8) circle (0.8) (1.6,8) circle (0.8);
  \draw[#3!80!black, line width=0.8pt] (0,11) -- (0,13); \fill[bkRed] (0,13.6) circle (0.9);
  \draw[#3!80!black, line width=1.6pt, line cap=round] (-5,1) -- (-8,-3) (5,1) -- (8,-3) (-2.5,-6) -- (-2.5,-10) (2.5,-6) -- (2.5,-10);
  \end{scope}}
% 路線の1行  \kviiaroute{y}{路線番号}{前の巻}{この巻の駅}{この先}
\newcommand\kviiaroute[5]{%
  \draw[bkRoute#2, line width=3pt, line cap=round] (10,#1) -- (132,#1);
  \node[inner sep=0pt] at (4,#1) {\bkrouteicon[7mm]{#2}};
  \fill[white] (24,#1) circle (2.2);\draw[bkRoute#2, line width=1.1pt] (24,#1) circle (2.2);
  \node[anchor=north, font=\sffamily\footnotesize, text=bkGray, inner sep=0pt] at (24,#1-3.8) {#3};
  \fill[white, rounded corners=1.5mm] (38,#1-4.4) rectangle (110,#1+4.4);
  \draw[bkRoute#2, line width=1.6pt, rounded corners=1.5mm] (38,#1-4.4) rectangle (110,#1+4.4);
  \node[font=\sffamily\footnotesize, text=bkInk] at (74,#1) {#4};
  \fill[white] (124,#1) circle (2.2);\draw[bkRoute#2, line width=1.1pt] (124,#1) circle (2.2);
  \node[anchor=north, font=\sffamily\footnotesize, text=bkGray, inner sep=0pt] at (124,#1-3.8) {#5};}
% 部の地図の1まい  \kviiapart{上のはしの y}{部の名前}{章}{絵}{中身}   はば132mm・高さ38mm
\newcommand\kviiapart[5]{%
  \draw[bkTheme!80!black, line width=1pt, fill=white, rounded corners=2.5mm] (0,#1) rectangle (132,#1-45);
  \node[anchor=north west, fill=bkTheme, text=bkOn, rounded corners=1.2mm, font=\sffamily\bfseries\small,
    inner xsep=2mm, inner ysep=1mm] at (3.5,#1-3.5) {#2};
  \node[anchor=north west, font=\sffamily\footnotesize, text=bkGray, inner sep=0pt] at (4.5,#1-11.5) {#3};
  \begin{scope}[shift={(23,#1-28)}]#4\end{scope}
  \node[anchor=north west, text width=78mm, align=left, font=\small, inner sep=0pt] at (51,#1-11) {#5};}
% 箇条書きの1項目  \kviiaitem{文}
\newcommand\kviiaitem[1]{\par\noindent\hangindent=1\zw\hangafter=1 ・#1\par}
% 会話の1行  \kviiatalk{y}{sora|mina|delta}{名前}{せりふ}
\newcommand\kviiatalk[4]{%
  \bkChar{#2}{(7,#1-6.5)}{12mm}
  \node[rectangle callout, callout absolute pointer={(15.5,#1)}, anchor=west, draw=bkTheme!80!black, line width=0.8pt,
    fill=white, rounded corners=2.5mm, inner xsep=3mm, inner ysep=1.8mm, text width=103mm, align=left,
    font=\small] at (21,#1) {#4};
  \node[anchor=north, font=\sffamily\footnotesize, text=bkGray, inner sep=0pt] at (7,#1-7.5) {#3};}
% 読み方のカード  \kviiacard{(x,y)}{アイコン}{色}{見出し}{中身}  はば64mm・高さ54mm
\newcommand\kviiacard[5]{%
  \node[anchor=north west, draw=#3, line width=1pt, fill=white, rounded corners=2.5mm,
    minimum width=64mm, minimum height=54mm, inner sep=0pt] at #1 {};
  \node[anchor=north west, inner sep=0pt] (ci) at ($#1+(3.5,-3)$) {\bkicon[1.5em]{#2}{#3}};
  \node[anchor=west, font=\sffamily\bfseries\small, text=#3!85!black, inner sep=0pt] at ([xshift=2mm]ci.east) {#4};
  \node[anchor=north west, text width=57mm, align=left, font=\small, inner sep=0pt] at ($#1+(3.5,-12)$) {#5};}
% 場面のカード  \kviiascene{(左上)}{章}{見出し}{絵}{説明}  はば64mm・高さ46mm
\newcommand\kviiascene[5]{%
  \draw[bkTheme!70!black, line width=0.9pt, fill=white, rounded corners=2.5mm] #1 rectangle ++(64,-46);
  \node[anchor=north west, fill=bkTheme, text=bkOn, rounded corners=1mm, font=\sffamily\bfseries\footnotesize, inner xsep=1.6mm, inner ysep=0.8mm] at ($#1+(3,-3)$) {#2};
  \node[anchor=north west, font=\sffamily\bfseries\small, text=\bkTT, inner sep=0pt] at ($#1+(17,-3.6)$) {#3};
  \begin{scope}[shift={($#1+(16,-27)$)}]#4\end{scope}
  \node[anchor=north west, text width=29mm, align=left, font=\footnotesize, inner sep=0pt] at ($#1+(33,-12)$) {#5};}
% 見出しのアイコン（路線：線と2つの駅）
\newcommand\bk@icon@kviiaroute{%
  \draw[bkOn, line width=0.13*\bk@iu, line cap=round] (-0.62,0) -- (0.62,0);
  \filldraw[fill=white, draw=bkOn, line width=0.09*\bk@iu] (-0.36,0) circle (0.2) (0.36,0) circle (0.2);}
\makeatother

\bkchapter{7}{0}{表のままで\\計算できたら？}{\bklevel{3}}{表のままで計算できたら？}

\noindent ゲームの画面の中で、キャラクターがくるっと回ったり、大きくなったり、ななめにかたむいたりします。車のゲームなら、ハンドルを切ると、道や建物の見え方がなめらかに変わっていきます。

\begin{bkfigure}
\begin{tikzpicture}[x=1mm, y=1mm]
  \filldraw[fill=bkInk!85, draw=bkInk, line width=0.8pt, rounded corners=4mm] (0,0) rectangle (112,48);
  \fill[bkBlue!8, rounded corners=2mm] (4,4) rectangle (108,44);
  \fill[bkGreen!25] (4,4) rectangle (108,12);
  \kviiachara{(26,22)}{0}{bkOrange}
  \kviiachara{(62,22)}{-30}{bkOrange}
  \kviiachara{(94,22)}{-90}{bkOrange}
  \draw[bkGray, line width=0.8pt, dashed, -{Stealth[length=2.4mm]}] (38,34) to[bend left=25] (52,34);
  \draw[bkGray, line width=0.8pt, dashed, -{Stealth[length=2.4mm]}] (72,34) to[bend left=25] (84,34);
  \node[anchor=west, text width=18mm, align=left, font=\sffamily\footnotesize, text=bkDeep] at (114,24) {くるっと\\回るとき、\\何を計算\\している？};
\end{tikzpicture}
\bkfigcaption{ゲームの画面。キャラクターが右回りに少しずつ回っていく。}
\end{bkfigure}

画面の絵は、たくさんの点の位置でできています。キャラクターが回ると、その点の位置がすべて、いっせいに変わります。しかも、画面は1秒に60回もかきかえられます。コンピュータは、どうやってそんなにたくさんの位置を、すばやく計算しているのでしょう。

\begin{hatena}
\par\addvspace{4pt}\noindent\begin{minipage}[c]{0.58\linewidth}
ゲームの画面で、キャラクターがくるっと回るとき、コンピュータは何を計算しているのかな？\par\vspace{3pt}
ロボットの腕や、空を飛ぶドローンの向きも、同じ方法で計算できるのかな？
\end{minipage}\hfill
\begin{minipage}[c]{0.38\linewidth}\centering
\begin{tikzpicture}[x=1mm, y=1mm]
  \kviiachara{(8,6)}{0}{bkOrange}
  \kviiachara{(36,6)}{-45}{bkOrange}
  \draw[bk arrow, line width=1pt] (17,10) -- node[above, font=\sffamily\bfseries\large, text=bkOrange!80!black, inner sep=1.5mm] {？} (27,10);
\end{tikzpicture}
\end{minipage}
\end{hatena}

答えのかぎは、この巻の主役「\textbf{行列}」です。行列は、数をたてと横に並べた表です。序章では、まず方眼の上の絵を回してみて、「表のまま計算する」とはどういうことかを見ます。そのあとで、これまでの巻の「表」をふり返り、この巻の地図を作ります。

\begin{omoidasou}
\bkrecall{1-13}{九九の表で、よこのならびを行、たてのならびを列という。「たて・よこの表でけいさんをまとめることを『行列』という。第7巻でまなぶ」。}
\bkrecall{4-8}{点 $(x,\ y)$ を原点のまわりに左回りに90°回すと、$(-y,\ x)$ になる。}
\bkrecall{3-6}{連立方程式の係数を表に並べて、表の行どうしをひいて解いた（係数の表は「行列の種」）。}
\bkrecall{6-14}{ばねでつながった2つの台車の式の係数を、表の形に並べた。}
\end{omoidasou}

\begin{yattemiyou}[方眼の家の絵を、90°回そう]
\noindent\begin{minipage}[c]{0.58\linewidth}
\raggedright
方眼に、家の絵をかきます。5つの頂点①〜⑤の座標は、下の表のとおりです（戸は右によせてかく）。この家を、原点Oのまわりに左回りに90°回します。
\begin{enumerate}[label={\bkcircled[bkGreen]{\arabic*}}, leftmargin=2.4\zw, itemsep=1pt, topsep=3pt]
\item 4-8 のきまり「$(x,\ y)$ は $(-y,\ x)$ になる」を使って、5つの頂点の新しい座標を表に書こう。
\item 新しい座標の点を方眼にかいて、順に結ぼう。家は回ったかな？
\end{enumerate}
\end{minipage}\hfill
\begin{minipage}[c]{0.38\linewidth}\centering
\begin{tikzpicture}[x=5mm, y=5mm]
  \kviiagrid{-4}{3}{-1}{4}
  \kviiahouse[bkGray]{1}{0}{0}{1}
  \foreach \x/\y/\n/\a in {0/0/1/south west, 2/0/2/south east, 2/2/3/west, 1/3/4/south, 0/2/5/east} {
    \fill[bkInk] (\x,\y) circle (0.12);}
  \node[kviia lab, anchor=north west] at (2.05,-0.05) {②};
  \node[kviia lab, anchor=north west] at (0.1,-0.05) {①};
  \node[kviia lab, anchor=west] at (2.15,2) {③};
  \node[kviia lab, anchor=south] at (1,3.15) {④};
  \node[kviia lab, anchor=south east] at (-0.1,2.05) {⑤};
\end{tikzpicture}
\end{minipage}
\begin{bkcenter}
\begin{bktable}{colspec={Q[l,m,wd=28mm] *{5}{Q[c,m,wd=14mm]}}, rowsep=1.5pt, row{1}={font=\sffamily\small}}
頂点 & ① & ② & ③ & ④ & ⑤ \\
もとの座標 & (0, 0) & (2, 0) & (2, 2) & (1, 3) & (0, 2) \\
回した座標 & ？ & ？ & ？ & ？ & ？ \\
\end{bktable}
\end{bkcenter}
\end{yattemiyou}

\noindent ①で、頂点②は $(2,\ 0)$ なので、$x=2$、$y=0$ です。$(-y,\ x)$ に入れると、新しい座標は $(0,\ 2)$ です。同じようにすると、次の表になります。

\begin{bkcenter}
\begin{bktable}{colspec={Q[l,m,wd=28mm] *{5}{Q[c,m,wd=15mm]}}, rowsep=1.5pt, row{1}={font=\sffamily\small}}
頂点 & ① & ② & ③ & ④ & ⑤ \\
もとの座標 & (0, 0) & (2, 0) & (2, 2) & (1, 3) & (0, 2) \\
回した座標 & (0, 0) & (0, 2) & ($-2$, 2) & ($-3$, 1) & ($-2$, 0) \\
\end{bktable}
\bktabcaption{家の5つの頂点を、左回りに90°回した座標}
\end{bkcenter}

\noindent\begin{minipage}[c]{0.52\linewidth}\raggedright
②で結ぶと、家は左に90°たおれた形になります（緑の家）。戸も、ちゃんと家といっしょに回っています。\par\vspace{3pt}
絵全体を回すのに、計算したのは5つの頂点だけです。5つの頂点に同じ計算をするだけで、絵全体が回ったのです。ゲームのキャラクターの絵は、何千もの点でできています。それでも、することは同じです。
\end{minipage}\hfill
\begin{minipage}[c]{0.44\linewidth}\centering
\begin{tikzpicture}[x=5mm, y=5mm]
  \kviiagrid{-4}{3}{-1}{4}
  \kviiahousezero
  \kviiahouse[bkGreen]{0}{-1}{1}{0}
  \node[kviia lab, anchor=south, text=bkGreen!50!black, fill=white] at (-2,2.3) {回した家};
  \node[kviia lab, anchor=south, text=bkGray, fill=white] at (1,3.3) {もとの家};
\end{tikzpicture}
\end{minipage}

%% ---------------- ① 4つの数のルール ----------------
\bksection{doguu}{4つの数で、絵を変える}

\bksubsection{1}{回すルールを、4つの数の表にする}

\label{bk:7-0:yotsu}やってみようで使ったルールを、もう一度くわしく書きます。新しい $x$ はもとの $-y$、新しい $y$ はもとの $x$ でした。これを、もとの $x$ と $y$ を両方使う形で書きなおします。
\[ \text{新しい}\ x=0\times x+(-1)\times y,\qquad \text{新しい}\ y=1\times x+0\times y \]
「$x$ に何をかけて、$y$ に何をかけて、たすか」を決める4つの数 0、$-1$、1、0 が出てきました。この4つの数を、式の並び方のまま、表にします。
\[ \begin{pmatrix} 0 & -1\\ 1 & 0 \end{pmatrix} \]
1行め（0、$-1$）が「新しい $x$ の作り方」、2行め（1、0）が「新しい $y$ の作り方」です。この4つの数の表を、どの頂点にも同じように使えば、絵全体が90°回ります。このように、数をたてと横に並べた表を\textbf{行列}といいます（7-1）。\bkidx{ぎょうれつ}{行列}

\bksubsection{2}{4つの数を変えると、絵の変わり方が変わる}

\label{bk:7-0:kawaru}4つの数を変えてみます。どの表も、「新しい $x$＝1行めの2つの数を $x$、$y$ にかけてたす」「新しい $y$＝2行めの2つの数を $x$、$y$ にかけてたす」というルールで使います。

\begin{bkfigure}
\begin{tikzpicture}[x=4mm, y=4mm]
  \begin{scope}
    \kviiagrid{-1}{5}{-1}{4}
    \kviiahousezero
    \kviiahouse[bkRed]{2}{0}{0}{1}
    \node[anchor=north] at (2,-1.6) {$\begin{pmatrix} 2 & 0\\ 0 & 1 \end{pmatrix}$};
    \node[kviia lab, anchor=north] at (2,-6.6) {横に2倍にのびる};
  \end{scope}
  \begin{scope}[xshift=38mm]
    \kviiagrid{-1}{5}{-1}{4}
    \kviiahousezero
    \kviiahouse[bkPurple]{1}{1}{0}{1}
    \node[anchor=north] at (2,-1.6) {$\begin{pmatrix} 1 & 1\\ 0 & 1 \end{pmatrix}$};
    \node[kviia lab, anchor=north] at (2,-6.6) {上ほど右へずれる};
  \end{scope}
  \begin{scope}[xshift=76mm]
    \kviiagrid{-1}{7}{-1}{8}
    \kviiahousezero
    \kviiahouse{2}{1}{1}{2}
    \node[anchor=north] at (3,-1.6) {$\begin{pmatrix} 2 & 1\\ 1 & 2 \end{pmatrix}$};
    \node[kviia lab, anchor=north] at (3,-6.6) {ななめにのびる};
  \end{scope}
\end{tikzpicture}
\bkfigcaption{4つの数を変えると、同じ家の絵が、ちがう形に変わる。灰色の点線が、もとの家。}
\label{bk:7-0:fig-kawaru}
\end{bkfigure}

図\ref{bk:7-0:fig-kawaru}の左の表では、新しい $x$ が $2x+0y=2x$、新しい $y$ が $0x+1y=y$ です。横だけが2倍になります。まん中の表では、新しい $x$ が $x+y$ で、上の点ほど右へずれます。右の表（主役の行列）は、家をななめ右上へ大きくのばしました。新しい $x$ が $2x+y$、新しい $y$ が $x+2y$ です。

図\ref{bk:7-0:fig-kawaru}の右の行列は、\textbf{この巻の主役}です。この巻では、この同じ行列を、章ごとにちがう目で見ていきます（⑧）。

\bksubsection{3}{コンピュータがしていること}

\label{bk:7-0:kotae}\noindent\textbf{\bkhatenaref の答え}\par
ゲームのコンピュータは、キャラクターの絵の点の1つ1つに、同じ「4つの数の表」を使って、新しい位置を計算しています。回す角が90°でないときは、4つの数が0や1ではなく、$\cos$ と $\sin$ の値（4-8）になります。たとえば30°回すなら、
\[ \begin{pmatrix} \cos30^\circ & -\sin30^\circ\\ \sin30^\circ & \cos30^\circ \end{pmatrix}\approx\begin{pmatrix} 0.866 & -0.5\\ 0.5 & 0.866 \end{pmatrix} \]
です。この表を\textbf{回転行列}といい、7-6 で作ります。画面が1秒に60回かきかわるたびに、少しずつ角を変えた回転行列で、何千、何万という点を計算しなおしているのです。

3次元のゲームでは、点の座標が $(x,\ y,\ z)$ の3つになり、3×3 や4×4 の表を使います（7-7）。ロボットの腕の関節や、ドローンの向きも、同じように行列で計算します。行列は、「たくさんの数を、同じルールで、いっせいに変える」ための道具なのです。

\begin{seikaku}
4つの数の表を点にかけるときの「1行めを使って新しい $x$、2行めを使って新しい $y$」という計算は、7-1 で「行列×ベクトル」として、ていねいに決めます。また、4つの数の表で表せるのは、次の3つがそろった変え方だけです。\textbf{まっすぐな線は、まっすぐな線のまま}（曲がらない）。\textbf{平行な線は、平行なまま}（等しい間かくで並んだ方眼の線は、変えたあとも等しい間かくで並ぶ）。\textbf{原点は動かない}。のばす・回す・ずらす（かたむける）・裏返すは、どれもこの仲間です。7-2 で、このきまりを「線形」の2つの式として確かめます。絵を曲げる変え方や、絵をそのまま横へ動かす「平行移動」は、4つの数だけではできません。平行移動は、7-7 で、数を1つふやす工夫をして行列で表します。
\end{seikaku}

%% ---------------- ② 九九の表 ----------------
\bksection{doguu}{九九の表から始まった「表」の道}

\bksubsection{4}{九九の表をもう一度}

\label{bk:7-0:kuku}第1巻 1-13 で、九九の表を作りました。$3\times4$ の答えは、左の3から横へ、上の4からたてへ進んで出会うマスにありました。横のならびを行、たてのならびを列とよびました。

\noindent\begin{minipage}[c]{0.4\linewidth}\centering
\begin{tikzpicture}[x=1mm, y=1mm]
  \foreach \i in {1,...,5} {\foreach \j in {1,...,5} {
    \pgfmathtruncatemacro\v{\i*\j}
    \ifnum\i=\j \def\cc{bkGreen!20}\else\def\cc{white}\fi
    \ifnum\i=3 \def\cc{bkOrange!20}\fi
    \ifnum\j=4 \def\cc{bkBlue!14}\fi
    \ifnum\i=3 \ifnum\j=4 \def\cc{bkGreen!35}\fi\fi
    \node[draw=bkInk!50, fill=\cc, line width=0.5pt, minimum size=7.5mm, inner sep=0pt, font=\small] at ({7.5*\j},{-7.5*\i}) {\v};}}
  \foreach \k in {1,...,5} {\node[kviia lab, text=bkGray] at ({7.5*\k},-1.6) {\k}; \node[kviia lab, text=bkGray] at (1.6,{-7.5*\k}) {\k};}
  \node[kviia lab, anchor=north, align=center] at (22.5,-42) {オレンジが3の行、青が4の列\\出会うマスが $3\times4=12$};
\end{tikzpicture}
\end{minipage}\hfill
\begin{minipage}[c]{0.56\linewidth}\raggedright
1-13 では、この表から4つのきまりを見つけました。そのうちの2つは、この巻でそのまま使います。\par\vspace{3pt}
・\textbf{ななめのきまり}：1×1、2×2、…は、左上から右下へのななめに並ぶ（緑）。この巻では、このななめの数を「対角成分」とよびます（7-1）。\par
・\textbf{入れかえても同じ}：$3\times4$ と $4\times3$ は同じ12。ななめをはさんで向かい合う数が等しい。この巻では「対称な行列」とよびます（7-1）。
\end{minipage}

\medskip
1-13 の終わりでは、時間わりの「水曜日の2時間め」が行と列の出会う所であること、テレビやゲームの画面が小さな点のたて・横の表であること、ロボットの腕の動きもコンピュータが数の表で計算していることを見ました。そして、デルタが「たて・よこの表でけいさんをまとめることを『行列』というよ。第7巻でまなぶ」と約束しました。この巻で、その約束をはたします。

\bksubsection{5}{表の路線をふり返る}

\label{bk:7-0:rosen5}この本の章は、7本の路線の駅になっています。「表」の路線（\bkrouteicon[0.9em]{5}）は、九九の表から始まって、この巻の 7-1 に着きます。

\begin{bkfigure}
\begin{tikzpicture}[x=1mm, y=1mm]
  \draw[bkRoute5, line width=3pt, line cap=round] (4,0) -- (128,0);
  \foreach \x/\c/\t/\s in {6/bkVol1/{1-13}/{九九の表}, 26/bkVol1/{1-17}/{表とグラフ}, 46/bkVol2/{2-18}/{表計算}, 66/bkVol3/{3-6}/{係数の表}, 86/bkVol6/{6-14}/{連立の式}, 106/bkVol7/{7-1}/{行列}, 126/bkVol9/{9-10}/{制御}} {
    \fill[white, draw=\c, line width=1.6pt] (\x,0) circle (2.6);
    \node[anchor=base, font=\sffamily\bfseries\footnotesize, text=\c!85!black] at (\x,-8) {\t};
    \node[anchor=base, font=\sffamily\footnotesize, text=bkInk] at (\x,-13.5) {\s};}
  \node[anchor=west, inner sep=0pt] at (0,9) {\bkrouteicon[1.4em]{5}\hspace{0.3em}{\sffamily\bfseries\small\color{bkRoute5!80!black}表の路線}};
\end{tikzpicture}
\bkfigcaption{表の路線。九九の表から、表計算、連立方程式の係数の表を通って、7-1 の行列に着く。この先は、第9巻の制御（状態方程式）へ続く。}
\end{bkfigure}

\begin{bkcenter}
\begin{bktable}{colspec={Q[c,m,wd=11mm] Q[l,m,wd=28mm] Q[l,m,wd=82mm]}, rowsep=1pt, cells={font=\small}, row{1}={font=\sffamily\bfseries\small}}
章 & 駅 & そこでしたこと \\
1-13 & 九九の表 & 行と列。ななめのきまり。入れかえても同じ \\
1-17 & 表とグラフ & 数を表に整理して、グラフにする \\
2-18 & 表計算 & セルに式を入れて、同じ計算をくり返す \\
3-6 & 連立方程式 & 係数だけを表に並べて、行どうしをひいて解く \\
6-14 & 連立微分方程式 & 2つの式の係数を表に並べる。和と差でほどける \\
\textbf{7-1} & \textbf{行列} & 表のまま、たし算・かけ算をする \\
\end{bktable}
\bktabcaption{表の路線の駅と、そこでしたこと}
\end{bkcenter}

3-6 では、「表の行どうしをひいたり、行に数をかけたりするのは、加減法そのもの。コンピュータはこのやり方で、文字が何千もある連立方程式を解いている。この数の表を、大学の数学では行列とよび、7-1 で学ぶ」と書きました。6-14 では、係数の表から「ほどける組み合わせ」を見つける方法が、第7巻の固有ベクトル（7-8）と対角化（7-9）だと予告しました。どちらも、この巻で回収します。

%% ---------------- ③ 行列が働く場面 ----------------
\bksection{butsuri}{行列が働く場面}

\bksubsection{6}{この巻で出会う4つの場面}

\label{bk:7-0:bamen}行列は、たくさんの数を「表のまま」まとめて計算する道具です。この巻では、次の4つの場面に何回も出会います。

\begin{bkfigure}
\begin{tikzpicture}[x=1mm, y=1mm]
  \kviiascene{(0,0)}{7-1}{買い物の合計}{%
    \foreach \i in {0,1} {\foreach \j in {0,1} {\draw[bkInk!50, line width=0.5pt, fill=bkOrange!15] ({-12+5*\j},{5-5*\i}) rectangle ++(4.6,-4.6);}}
    \node[font=\small] at (0,0) {$\times$};
    \foreach \i in {0,1} {\foreach \j in {0,1} {\draw[bkInk!50, line width=0.5pt, fill=bkBlue!12] ({3+5*\j},{5-5*\i}) rectangle ++(4.6,-4.6);}}}%
    {個数の表×値段の表で、何人分・何店分の合計を一度に出す。}
  \kviiascene{(68,0)}{7-2}{絵をのばす・回す}{%
    \begin{scope}[shift={(-7,-10)}, x=2.4mm, y=2.4mm]
      \kviiahousezero
      \kviiahouse{2}{1}{1}{2}
    \end{scope}}%
    {方眼の絵の点に行列をかけて、形を変える（7-2・7-6）。}
  \kviiascene{(0,-50)}{7-7}{ロボットとドローン}{%
    \fill[bkGray!35] (-12,-9) rectangle (2,-7);
    \draw[bkOrange!80!black, line width=2.6pt, line cap=round] (-5,-7) -- ++(70:10) coordinate (E);
    \draw[bkBlue!80!black, line width=2.6pt, line cap=round] (E) -- ++(10:9);
    \filldraw[fill=white, draw=bkInk, line width=0.7pt] (-5,-7) circle (1.2) (E) circle (1.2);
    \begin{scope}[shift={(4,7)}, rotate=-12]
      \draw[bkInk, line width=1pt] (-6,0) -- (6,0);
      \fill[bkInk!70] (-2,-1.2) rectangle (2,1.2);
      \draw[bkInk, line width=0.8pt] (-8,1.2) -- (-4,1.2) (4,1.2) -- (8,1.2);
    \end{scope}}%
    {関節の角や機体の向き（3次元の回転）を、行列のかけ算で計算する。}
  \kviiascene{(68,-50)}{7-8}{つながったゆれ}{%
    \draw[bkInk, line width=0.9pt] (-14,-6) -- (-14,6);
    \draw[bkInk, line width=0.7pt, decorate, decoration={zigzag, amplitude=0.9mm, segment length=1.6mm}] (-14,0) -- (-9,0);
    \filldraw[fill=bkOrange!30, draw=bkOrange!80!black, line width=0.7pt] (-9,-3.5) rectangle (-3,3.5);
    \draw[bkInk, line width=0.7pt, decorate, decoration={zigzag, amplitude=0.9mm, segment length=1.6mm}] (-3,0) -- (2,0);
    \filldraw[fill=bkBlue!25, draw=bkBlue!80!black, line width=0.7pt] (2,-3.5) rectangle (8,3.5);
    \draw[bkInk, line width=0.7pt, decorate, decoration={zigzag, amplitude=0.9mm, segment length=1.6mm}] (8,0) -- (12,0);
    \draw[bkInk, line width=0.9pt] (12,-6) -- (12,6);}%
    {ばねでつながった台車や建物のゆれ方（振動モード）を、行列から見つける。}
\end{tikzpicture}
\bkfigcaption{この巻で行列が働く4つの場面}
\end{bkfigure}

どの場面でも、数が2つ、3つ、…と組になっていて、それを同じルールでいっせいに計算します。1つずつの式で書くと何行にもなる計算が、行列を使うと、$A\vec x$ のような短い式で書けます。そして、表計算や Python で、何百、何千という数の計算を一度にできます。

%% ---------------- ④ この巻の地図 ----------------
\bksection{goal}{この巻の地図}

\bksubsection{7}{3つの部}

\label{bk:7-0:chizu}第7巻は、3つの部に分かれています。7-1〜7-6 は高校の数学で学ぶことがある内容、7-7〜7-10 は大学（工学部）で学ぶ内容です。大学の内容の章は、「目でわかる」「手でわかる」「式でわかる」の3段に分けてあり、「ここまでで十分」の旗の所まで読めば、その章のゴールです（第6巻と同じ）。

\begin{bkfigure}
\begin{tikzpicture}[x=1mm, y=1mm]
  \kviiapart{0}{第1部\enspace 行列の計算}{7-1〜7-5}%
    {\matrix[kviia mat, nodes={minimum width=6mm, minimum height=5mm, font=\small}] at (0,0) {2 \& 1\\ 1 \& 2\\};}
    {\kviiaitem{数の表（行列）のたし算・かけ算（7-1）。}\kviiaitem{行列を「形を変える機械」として見る。かけ算は変換をつなぐこと（7-2・7-3）。}\kviiaitem{面積の倍率（行列式）と、元にもどす行列（逆行列）。連立方程式を解く（7-4・7-5）。}}
  \kviiapart{-49}{第2部\enspace 回転と座標の変換}{7-6〜7-7}%
    {\draw[bkInk, line width=0.5pt] (-12,0) -- (12,0) (0,-10) -- (0,10);
     \draw[bkInk!70, line width=0.6pt] (0,0) circle (8);
     \draw[-{Stealth[length=2.4mm]}, bkRed, line width=1.2pt] (0,0) -- (10:8);
     \draw[-{Stealth[length=2.4mm]}, bkGreen!60!black, line width=1.2pt] (0,0) -- (55:8);
     \draw[bkInk, line width=0.5pt, -{Stealth[length=1.6mm]}] (10:4) arc (10:55:4);}
    {\kviiaitem{回転行列で、点や絵をくるっと回す（7-6）。}\kviiaitem{3次元の回転。ドローンの向き（ロール・ピッチ・ヨー）と、ロボットの腕の位置（同次変換行列）（7-7）。}}
  \kviiapart{-98}{第3部\enspace 固有値と行列の使いみち}{7-8〜7-11}%
    {\draw[bkInk, line width=0.5pt] (-12,0) -- (12,0) (0,-10) -- (0,10);
     \draw[-{Stealth[length=2.4mm]}, bkOrange!80!black, line width=1.2pt] (0,0) -- (45:10);
     \draw[-{Stealth[length=2.4mm]}, bkBlue!80!black, line width=1.2pt] (0,0) -- (-45:5);
     \draw[bkGray, line width=0.6pt, dashed] (225:9) -- (45:12);}
    {\kviiaitem{行列をかけても向きが変わらない矢印（固有ベクトル）（7-8）。}\kviiaitem{その矢印を物差しにすると、計算がほどける（7-9）。}\kviiaitem{何回もかける（7-10）。工学での使いみち（7-11）。}}
\end{tikzpicture}
\bkfigcaption{第7巻の3つの部と、それぞれで学ぶこと}
\end{bkfigure}

\bksubsection{8}{主役の行列の、6つの顔}

\label{bk:7-0:shuyaku}この巻には、主役の行列がいます。②で家の絵をななめにのばした行列です。
\[ A=\begin{pmatrix} 2 & 1\\ 1 & 2 \end{pmatrix} \]
この巻では、同じこの行列を、章ごとにちがう目で見ていきます。同じ人でも、学校・家・部活で見せる顔がちがうように、1つの行列にも、いくつもの顔があるのです。

\begin{bkcenter}
\begin{bktable}{colspec={Q[c,m,wd=11mm] Q[l,m,wd=42mm] Q[l,m,wd=70mm]}, rowsep=1.2pt, cells={font=\small}, row{1}={font=\sffamily\bfseries\small}}
章 & 見る目 & 主役の行列 $A$ の顔 \\
7-2 & 形を変える機械 & 方眼の家を、ななめ右上へのばす \\
7-4 & 面積が何倍になるか & 方眼の正方形1つが、面積3の平行四辺形になる \\
7-5 & 元にもどす機械 & のびた家を、元の家にもどす行列がある \\
7-8 & 向きの変わらない矢印 & ななめ45°の矢印は3倍にのびるだけ、その直角の向きの矢印は1倍（そのまま） \\
7-9 & のびるだけに分ける & ななめの物差しで見ると「3倍と1倍にのばすだけ」 \\
7-10 & 何回もかける & 何回もかけると、点はななめ45°の向きにそろっていく \\
\end{bktable}
\bktabcaption{主役の行列の6つの顔（くわしくは、それぞれの章で）}
\end{bkcenter}

いまは、表の中のことばの意味がわからなくてもだいじょうぶです。章を読み進めて、この表にもどってくると、1つの行列が少しずつ立体的に見えてくるはずです。

%% ---------------- ⑤ 道具 ----------------
\bksection{doguu}{表計算と Python}

\bksubsection{9}{手で計算して、コンピュータで確かめる}

\label{bk:7-0:dougu}行列の計算は、1つ1つはかけ算とたし算だけです。手で計算できます。でも、行列が大きくなると、計算の数がとても多くなります。そこで、この巻では、手で小さな例を計算したら、表計算（2-18）と Python（2-19）で同じ答えになるか確かめる、という進め方をします。

\noindent\begin{minipage}[t]{0.47\linewidth}\vspace{0pt}
\textbf{表計算}：行列のかけ算は MMULT という関数でできます（7-1）。\par\vspace{3pt}
\centering
\begin{tikzpicture}[x=1mm, y=1mm]
  \node[draw=bkLine, line width=0.6pt, minimum height=5.5mm, minimum width=10mm, font=\sffamily\footnotesize] (c) at (0,0) {D2};
  \node[anchor=west, draw=bkLine, line width=0.6pt, minimum height=5.5mm, minimum width=38mm, font=\ttfamily\footnotesize] at ([xshift=1mm]c.east) {=MMULT(A2:B3,C2:C3)};
\end{tikzpicture}
\end{minipage}\hfill
\begin{minipage}[t]{0.5\linewidth}\vspace{0pt}
\textbf{Python}：numpy という道具を使い、かけ算は \texttt{@} で書きます（7-1）。
\begin{bkcode}[Python]
import numpy as np
A = np.array([[2, 1],
              [1, 2]])
print(A @ np.array([1, 0]))
\end{bkcode}
\begin{bkcodeout}[実行結果]
[2 1]
\end{bkcodeout}
\end{minipage}

\begin{column}{「行列」という名前}
英語では、行列を matrix（メイトリクス）といいます。もとは「何かを生み出すもと」「型」という意味のことばです。数学のことばとして使ったのは、イギリスの数学者シルベスターで、1850年のことです。数の表から、いろいろな行列式（7-4）を生み出す「もと」という意味でした。そのあと、同じイギリスのケイリーが、1858年の論文で、行列どうしのたし算やかけ算のきまりをまとめました。日本語の「行列」は、「行」と「列」でできた表、という意味の訳語です。

ただし、数を表に並べて計算する考えは、もっとずっと古くからありました。2000年ほど前の中国の本『九章算術』には、連立方程式の係数を表に並べて解く方法が書かれています（7-1 のコラム）。
\end{column}

%% ---------------- ⑥ 新しい記号 ----------------
\bksection{card}{この巻で出会う新しい記号}

\label{bk:7-0:kigo}第7巻では、次のような新しい記号に出会います。いまは意味がわからなくてもだいじょうぶです。どの記号も、使う前に、数の例と図でていねいに説明します（記号カードは巻末にも並べてあります）。

\begin{bkcenter}
\begin{bktable}{colspec={Q[c,m,wd=34mm] Q[l,m,wd=30mm] Q[l,m,wd=48mm] Q[c,m,wd=9mm]}, rowsep=0.3pt, colsep=3pt, cells={font=\small}, row{1}={font=\sffamily\bfseries\small}}
記号 & 読み方 & ひとことで言うと & 章 \\
$A$、$a_{ij}$ & エー、エー・アイ・ジェイ & 行列と、その $i$ 行め $j$ 列めの成分 & 7-1 \\
$A^{\mathsf T}$ & エー・トランスポーズ & 行と列を入れかえた行列（転置） & 7-1 \\
$A\vec x$ & エー・エックス & 行列×ベクトル。形を変える機械 & 7-2 \\
$AB$、$I$ & エー・ビー、アイ & 行列の積、単位行列 & 7-3 \\
$\det A$ & デターミナント・エー & 面積が何倍になるか（行列式） & 7-4 \\
$A^{-1}$ & エー・インバース & 元にもどす行列（逆行列） & 7-5 \\
$R(\theta)$ & アール・シータ & $\theta$ だけ回す行列（回転行列） & 7-6 \\
$R_x$、$R_y$、$R_z$ & アール・エックスなど & 3次元の軸のまわりの回転 & 7-7 \\
$\lambda$、$\vec v$ & ラムダ、ブイ & 固有値と固有ベクトル & 7-8 \\
$P^{-1}AP=D$ & （対角化） & ななめの物差しで見ると、のびるだけ & 7-9 \\
$A^n$ & エーのエヌ乗 & 行列を $n$ 回かける & 7-10 \\
\end{bktable}
\bktabcaption{この巻で出会う記号（くわしくはそれぞれの章で）}
\end{bkcenter}

この巻では、行列は細い字の大文字（$A$、$R$、$J$、$K$ など）で、ベクトルは第4巻と同じく矢印つきの文字（$\vec x$）で書きます。ベクトルの成分は、この巻ではたてに並べます（7-1）。大学の本では、行列やベクトルを太字で書くことが多いので、7-11 で見くらべます。

%% ---------------- ⑦ 路線 ----------------
\bksection{kviiaroute}{この巻が通る路線と乗りかえ駅}

\label{bk:7-0:rosen}第7巻を通る路線は3本です。いちばん長く通るのは「表」の路線（\bkrouteicon[0.9em]{5}）で、7-6 は「回転と波」、7-7 は「矢印」の路線の駅です。

\begin{bkfigure}
\begin{tikzpicture}[x=1mm, y=1mm]
  \kviiaroute{28}{4}{6-10}{{\bfseries 7-6}}{9-1}
  \kviiaroute{14}{5}{3-6}{{\bfseries 7-1}・7-2〜7-5・7-8〜7-11}{9-10}
  \kviiaroute{0}{6}{4-10}{{\bfseries 7-7}}{8-5}
  \node[font=\sffamily\bfseries\small, text=bkGray] at (24,38) {前の巻};
  \node[font=\sffamily\bfseries\small, text=bkDeep] at (74,38) {この巻（第7巻）};
  \node[font=\sffamily\bfseries\small, text=bkGray] at (124,38) {この先};
\end{tikzpicture}
\bkfigcaption{第7巻を通る3本の路線。四角の中がこの巻の駅（太字は大きな駅）。}
\end{bkfigure}

\begin{bkcenter}
\begin{bktable}{colspec={Q[c,m,wd=7mm] Q[l,m,wd=20mm] Q[l,m,wd=92mm]}, rowsep=0.8pt, cells={font=\small}, row{1}={font=\sffamily\bfseries\small}}
 & 路線 & 大きな駅のつながり \\
\bkrouteicon{4} & 回転と波 & 三角関数（4-8）→ オイラーの公式（6-10）→ \textbf{回転行列}（7-6） \\
\bkrouteicon{5} & 表 & 表計算（2-18）→ 連立方程式（3-6）→ \textbf{行列}（7-1） \\
\bkrouteicon{6} & 矢印 & 座標（3-2）→ ベクトル（4-10）→ \textbf{3次元の回転}（7-7） \\
\end{bktable}
\bktabcaption{3本の路線の、この巻の大きな駅}
\end{bkcenter}

回転と波の路線では、6-9 で「複素数のかけ算は回転」と学びました。7-6 では、同じ回転を、複素数ではなく4つの数の表（回転行列）で表します。矢印の路線では、第4巻の2次元・3次元のベクトルが、7-7 で3次元の回転に出会い、第8巻のベクトル場（8-5）へ続きます。

%% ---------------- この巻の読み方 ----------------
\Needspace{22\baselineskip}\bksection{hon}{この巻の読み方}

\label{bk:7-0:yomikata}
\begin{bkfigure}
\begin{tikzpicture}[x=1mm, y=1mm]
  \kviiatalk{0}{sora}{ソラ}{第7巻は、大学の線形代数の内容なんだね。記号が多くて、むずかしそう。}
  \kviiatalk{-22}{mina}{ミナ}{でも、使う計算は、かけ算とたし算だけなんだって。九九の表と同じように、行と列を指でなぞればいいみたい。}
  \kviiatalk{-46}{delta}{デルタ}{そのとおり。1回読んで分からなくても、だいじょうぶ。小さな数の表で手を動かして、表計算や Python で確かめれば、かならずわかるよ。}
\end{tikzpicture}
\end{bkfigure}

\begin{bkfigure}
\begin{tikzpicture}[x=1mm, y=1mm]
  \kviiacard{(0,0)}{shiki}{bkNavy}{①\enspace 成分で書いてから記号に}
    {記号 $A\vec x$ だけで考えず、まず成分で1つずつ書きます。\par\vspace{2pt}
     \kviiaitem{2×2 の行列×ベクトルなら、$ax+by$ と $cx+dy$ を書く。}\kviiaitem{記号の式のあとには、2×2 の数で1回確かめる。}}
  \kviiacard{(68,0)}{me}{bkGreen}{②\enspace 図と数の2通りで}
    {初めての考えは、方眼の図と数の表の2通りで見せます。\par\vspace{2pt}
     \kviiaitem{家の絵が、どう変わるかを見る。}\kviiaitem{同じことを、数の計算で確かめる。}}
  \kviiacard{(0,-57)}{te}{bkOrange}{③\enspace 表計算と Python}
    {手で計算した答えを、コンピュータで確かめます。\par\vspace{2pt}
     \kviiaitem{表計算は MMULT など、Python は numpy の \texttt{@}。}\kviiaitem{型は 7-1 で決めて、全部の章で同じに使う。}}
  \kviiacard{(68,-57)}{goal}{bkDeep}{④\enspace ★5の章は3段}
    {大学の内容の章は、3つの段に分けてあります。\par\vspace{2pt}
     \kviiaitem{「目でわかる」「手でわかる」「式でわかる」。}\kviiaitem{「ここまでで十分」の旗まで読めば、その章のゴール。}}
\end{tikzpicture}
\end{bkfigure}
